\clearpage
\section{Compatible serial renaming and structural safety}
\label{sec:copying-serial-renaming}
This milestone closes the compatibility argument for renaming the original
serial assembly. Suppose $G:I\to J$ and $H:J\to K$ are replaced by boundary
isomorphic graphs $G':I'\to J'$ and $H':J'\to K'$. The two isomorphisms must
send each shared position of $J$ to the same position of $J'$. Under precisely
this condition, the original two-hole expansion admits a boundary isomorphism.
All ordered incidences and all recursive edge labels are preserved. No identity
wire, contraction, weakening, or box is removed.

The result supplies a composition step needed by Figure~16's composite clauses;
it does not by itself finish their renaming induction. We also prove that typed
boundary renaming preserves stability, binding-freedom and exact structural
paths. These are not machine-execution traces or a proof of Proposition~12.2.

\subsection{Extracting the three carrier actions}
\termdef{copying-boundary-components}
A boundary isomorphism preserves the predicates ``is an input'' and ``is an
output''. Consequently an internal port must map to an internal port; either
other result would contradict role preservation. Pattern matching on the port
image therefore extracts a map on internal vertices, and the equation saying
that this map reconstructs the original port image follows by the same cases.
The input and output maps are extracted similarly, using the positive role
condition to exclude the other two constructors.

\lean{Copying/BoundaryComponents}{8}{12}{BoundaryComponents.internalMap}
\lean{Copying/BoundaryComponents}{14}{20}{BoundaryComponents.map_internal}
\lean{Copying/BoundaryComponents}{22}{26}{BoundaryComponents.inputMap}
\lean{Copying/BoundaryComponents}{28}{33}{BoundaryComponents.map_input}
\lean{Copying/BoundaryComponents}{35}{39}{BoundaryComponents.outputMap}
\lean{Copying/BoundaryComponents}{41}{46}{BoundaryComponents.map_output}

Apply those constructions to the inverse isomorphism to obtain inverse
maps. Embed an alleged inverse identity into the corresponding port constructor,
rewrite both extracted maps as the original port map, and use the original
bijection's inverse identity. Injectivity of each constructor gives the desired
identity on the carrier. This constructs bijections on internal vertices,
input indices and output indices separately. Finally, specialising the original
vertex-label law to each constructor proves type preservation for all three
bijections. Thus compatibility below compares actual typed positions, not merely
interface cardinalities.

\lean{Copying/BoundaryComponents}{48}{58}{BoundaryComponents.internals}
\lean{Copying/BoundaryComponents}{60}{70}{BoundaryComponents.inputs}
\lean{Copying/BoundaryComponents}{72}{82}{BoundaryComponents.outputs}
\lean{Copying/BoundaryComponents}{84}{88}{BoundaryComponents.internal_label}
\lean{Copying/BoundaryComponents}{90}{93}{BoundaryComponents.input_label}
\lean{Copying/BoundaryComponents}{95}{98}{BoundaryComponents.output_label}

\Needspace{12\baselineskip}
\subsection{The exact gluing condition}
\termdef{copying-serial-compatible}
Let $f$ and $g$ be the two component isomorphisms. Write $f_O$ for the output
index bijection of $f$, and $g_I$ for the input index bijection of $g$.
Their compatibility is
\[
                 \forall j\in\operatorname{Fin}|J|,\qquad f_O(j)=g_I(j).
\]
It is inherited by inverses: apply the injective input map to the proposed
identity of inverse images, then use compatibility and both right-inverse laws.
It is inherited by composition because extracting an input or output map
commutes with composing port bijections. Equality of the two composite actions
then follows by the two given compatibility equations.

\lean{Copying/BoundarySerial}{13}{13}{BoundarySerial.Compatible}
\lean{Copying/BoundarySerial}{16}{20}{BoundarySerial.compatible_symm}
\lean{Copying/BoundarySerial}{23}{39}{BoundarySerial.compatible_trans}

The existing serial graph has internal vertices comprising shared-interface
positions, tagged left internal vertices, and tagged right internal vertices.
Its boundary consists of the left inputs and right outputs. Map these five
classes by $f_O$, the internal maps of $f$ and $g$, $f_I$, and $g_O$, respectively.
The inverse uses the five inverse bijections. Checking each constructor proves
both inverse laws. The edges retain their left/right tags and are mapped by the
appropriate component edge bijection.

\lean{Copying/BoundarySerial}{41}{69}{BoundarySerial.ports}
\lean{Copying/BoundarySerial}{138}{144}{BoundarySerial.edges}

The two insertion maps are the actual maps used by the existing two-hole
expansion. For the left insertion, the new port map commutes on internal,
input, and output ports by its definition. For the right insertion, the
internal and output cases are again direct; its input case is exactly the
compatibility equation, since that input has become a shared internal vertex.
Conversely, specialise commutation with the right insertion to each right input.
Injectivity of the internal-port and sum constructors recovers compatibility.
Thus agreement is necessary as well as sufficient for this explicit map to
respect both insertions. The converse does not claim that every possible
isomorphism between the assembled graphs must arise from these components.

\lean{Copying/BoundarySerial}{71}{73}{BoundarySerial.leftInsert}
\lean{Copying/BoundarySerial}{75}{77}{BoundarySerial.rightInsert}
\lean{Copying/BoundarySerial}{79}{98}{BoundarySerial.left_commutes}
\lean{Copying/BoundarySerial}{100}{120}{BoundarySerial.right_commutes}
\lean{Copying/BoundarySerial}{124}{136}{BoundarySerial.compatible_iff_right_commutes}

To prove the assembled isomorphism, type preservation is checked on the
same five port classes. Boundary roles follow immediately from their constructors.
An edge retains its component tag, so its recursive label relation is the
component's label relation. Its ordered source list is the component source list
mapped by its insertion. Use insertion commutation, the component's source law,
and composition of list maps. The target proof is identical with target lists.
Finally, evaluation on external inputs and outputs gives their exact actions:
$f_I$ on inputs, $g_O$ on outputs. All these proofs work at every allowance,
including zero, without requiring any interface to be inhabited.

\lean{Copying/BoundarySerial}{147}{186}{BoundarySerial.iso}
\lean{Copying/BoundarySerial}{189}{191}{BoundarySerial.iso_input}
\lean{Copying/BoundarySerial}{194}{196}{BoundarySerial.iso_output}

\subsection{Paths, nested binding restrictions, and stability}
For a boundary isomorphism $f:G\to H$, membership of the image of a port in
an image edge's source list is equivalent to original source membership.
Rewrite the source law and use injectivity of the port map to reflect membership
in the mapped list. Targets are identical. The existing adjacency equivalence
then transports ordered edge paths and exact-length adjacency traces. A vertex
step consists of a common edge witnessing source and target membership; transport
that edge forward, or recover it with the inverse edge bijection. This proves
vertex-step equivalence, hence exact vertex-trace and reachability equivalence.
The length preserved here counts graph incidences, not UAM transition costs.

\lean{Copying/BoundarySafety}{12}{19}{BoundarySafety.source_mem_iff}
\lean{Copying/BoundarySafety}{21}{28}{BoundarySafety.target_mem_iff}
\lean{Copying/BoundarySafety}{31}{33}{BoundarySafety.path_iff}
\lean{Copying/BoundarySafety}{36}{38}{BoundarySafety.edge_trace_iff}
\lean{Copying/BoundarySafety}{40}{48}{BoundarySafety.vertex_step_iff}
\lean{Copying/BoundarySafety}{50}{52}{BoundarySafety.vertex_trace_iff}
\lean{Copying/BoundarySafety}{54}{56}{BoundarySafety.vertex_reach_iff}

Every recursive primitive-label property is invariant: move each edge
back through the edge bijection and apply the existing recursive label-equivalence
law. Focus-freedom and hole-freedom are special cases. For shallow binding-freedom,
pull forbidden ancestry back using adjacency preservation, preserve binding and
hole labels, and transport the explicit source witness by source membership.
The proof retains the source requirement; it does not strengthen Lemma~E.1 at
zero-input holes. For a nested box, the existing recursive label relation provides
a boundary-fixing isomorphism of its body, to which the already proved nested
binding-path law applies. Combine this with the shallow result and focus-freedom,
and apply the inverse isomorphism, to obtain full binding-freedom in both directions.

\lean{Copying/BoundarySafety}{59}{65}{BoundarySafety.allLabels_iff}
\lean{Copying/BoundarySafety}{67}{67}{BoundarySafety.focusFree_iff}
\lean{Copying/BoundarySafety}{69}{69}{BoundarySafety.holeFree_iff}
\lean{Copying/BoundarySafety}{72}{77}{BoundarySafety.shallow_bindingFree}
\lean{Copying/BoundarySafety}{79}{98}{BoundarySafety.bindingPaths}
\lean{Copying/BoundarySafety}{101}{103}{BoundarySafety.bindingFree_iff}

Stability is also invariant when atom or thunk outputs are permuted.
Transport the permitted flat labels through the edge bijection and apply the
previously established characterisation of stability by these labels for the
example alphabet. That characterisation supplies the actual output and
all-vertex reachability conditions; neither condition is dropped from the
definition. Inversion proves reflection. This does not assert that arbitrary
execution steps preserve stability, nor does it alter the copyability predicate.

\lean{Copying/BoundarySafety}{106}{115}{BoundarySafety.stable_iff}

\subsection{Recovering the boundary-fixing interface}
If an isomorphism's extracted input and output maps fix every position, its
internal bijection and edge bijection form an existing recursive isomorphism.
For ordered incidences, compare the all-port map with the extension of the
internal map: they agree internally by extraction and on the boundary by the
two fixing hypotheses. Applying this to compatible serial assembly needs only
the left external inputs and right external outputs to be fixed. Its internal
shared-interface action may be nontrivial. The exact external action laws proved
above supply the two fixing hypotheses. This is a proved bridge to the original
observation infrastructure; the separate execution-invariance assumptions of
that infrastructure remain to be discharged by the concrete machine.

\lean{Copying/BoundarySafety}{119}{140}{BoundaryIso.toDeepIso}
\lean{Copying/BoundarySafety}{148}{161}{BoundarySerial.toDeepIso}

\subsection{Concrete distributor compositions and mismatch rejection}
\termdef{copying-distributor-serial-renaming}
The input-action transport lemma eliminates the four interface equalities and
then uses the original action equation. We use it to expose $D_{k,1}$ at interface
$k\to k$, by the equality $k\cdot1=k$ only. The graph retains its full forest of
weakening and contraction edges. Its isomorphism is the original distributor
resource automorphism, transported through that equality.

\lean{Copying/DistributorSerialRenaming}{9}{14}{DistributorSerialRenaming.castInterfaces_input}
\lean{Copying/DistributorSerialRenaming}{20}{21}{DistributorSerialRenaming.single}
\lean{Copying/DistributorSerialRenaming}{23}{25}{DistributorSerialRenaming.singleIso}

For a resource bijection $\pi$, the one-use forest sends input $i$ and
output $i$ to $\pi(i)$. The input position formula $i+k\cdot0=i$ and the previous
distributor output law prove these claims. Injectivity of the port constructors
then gives the same formula for the extracted input and output index maps.

\lean{Copying/DistributorSerialRenaming}{27}{35}{DistributorSerialRenaming.singleIso_input}
\lean{Copying/DistributorSerialRenaming}{37}{40}{DistributorSerialRenaming.singleIso_output}
\lean{Copying/DistributorSerialRenaming}{43}{50}{DistributorSerialRenaming.outputIndex_eq}
\lean{Copying/DistributorSerialRenaming}{53}{60}{DistributorSerialRenaming.singleInputIndex_eq}

Consequently the original $D_{k,m}$ followed by the original $D_{k,1}$
admits the assembled automorphism when both forests use $\pi$. This is an
unconditional family for every finite $k,m$ and every allowance; matching is
proved, not supplied as an application premise. Conversely, if the first forest
uses $\pi$ and the second uses $\tau$, compatibility holds exactly when the two
resource functions agree. Substitute the two explicit index formulas and compare
their finite-index values. The reverse direction substitutes equality of the
functions. This validates the gluing condition on actual sharing graphs rather
than only on an abstract interface.

\lean{Copying/DistributorSerialRenaming}{63}{66}{DistributorSerialRenaming.single_compatible}
\lean{Copying/DistributorSerialRenaming}{70}{73}{DistributorSerialRenaming.serial_singleIso}
\lean{Copying/DistributorSerialRenaming}{76}{86}{DistributorSerialRenaming.single_compatible_iff}

The regression family uses the involution $i\mapsto1-i$ on two resources.
Finite arithmetic proves its two inverse identities. Apply the distributor
composition theorem to two-use atom forests at arbitrary depth. A swap on one
side and identity on the other cannot satisfy compatibility: the characterisation
would force their functions to agree, and evaluation at zero gives $1=0$.
Separate instances cover no resources and no uses, so the theorem includes empty
forests and pure weakening. Finally, combining the generic serial isomorphism
with recursive binding-freedom invariance proves the safety equivalence of the
two actual serial assemblies, with the compatibility condition explicit.

\lean{Copying/BoundaryRenamingExamples}{8}{12}{BoundaryRenamingExamples.swapTwo}
\lean{Copying/BoundaryRenamingExamples}{15}{15}{BoundaryRenamingExamples.twoAtomUses}
\lean{Copying/BoundaryRenamingExamples}{18}{25}{BoundaryRenamingExamples.unmatched_swap}
\lean{Copying/BoundaryRenamingExamples}{28}{29}{BoundaryRenamingExamples.emptyForest}
\lean{Copying/BoundaryRenamingExamples}{32}{33}{BoundaryRenamingExamples.weakeningForest}
\lean{Copying/BoundaryRenamingExamples}{37}{42}{BoundaryRenamingExamples.serial_bindingFree_iff}

\paragraph{Remaining obligations and effort estimate.}
The completed subtask is typed serial gluing, including its exact compatibility
criterion, structural safety transport, and concrete distributor instances.
Whole-term renaming still requires the composite clauses and lambda boxes;
non-bijective renaming and graph substitution remain open. In particular, the
current recursive label relation fixes a box body's ordered boundary: a
permutation outside a box does not automatically justify permuting its body.
The induction must address that interface explicitly while retaining Figure~16.
As of 25 October 2026, estimated completion of the requested Proposition~12.2
formalisation is roughly \textbf{40--45\% by mathematical effort}, with a broad
plausible range of \textbf{30--50\%}. This is a judgement about remaining proof
work, not a count of declarations. The concrete copying machine and eleven
template families still account for a large part of the remaining effort.
